\documentstyle[11pt]{article}
\oddsidemargin -1mm
\evensidemargin -1mm
\topmargin 0mm
\headheight 0pt
\headsep 0pt
\footheight 12pt
\footskip 30pt
\textheight 234mm
\textwidth 160mm
\columnsep 10pt
\columnseprule 0pt
\topsep 1pt plus 2pt minus 4pt
\itemsep 1pt plus 2pt minus 1pt
\marginparwidth 0pt
\oddsidemargin -.5cm
\evensidemargin -.5cm
\marginparsep 0pt
\topmargin -.5cm
\textwidth 17cm
\textheight 23cm
\sloppy
\font\eirm=cmr8

\def \ce{\centering}
\def \N{{\rm I\kern -2.1pt N\hskip 1pt}}
\def \R{{\rm I\kern -2.2pt R\hskip 1pt}}
\def \Z{{\rm Z\!\!Z}}
\def \Q{\hskip2pt {\rm 0\kern -7pt Q}}
\def \C{{\rm C\kern -4.6pt
     \vrule height 6.8pt width 0.3pt depth -0.5pt}\hskip4pt}
\def \P{{\rm I\kern -2.2pt P\hskip 1pt}}
\def \K{{\rm I\kern -2.2pt K\hskip 1pt}}
\def \E{{\rm I\kern -2.2pt E\hskip 1pt}}
\def \H{{\rm I\kern -2.2pt H\hskip 1pt}}
\def \L{{\rm I\kern -2.2pt L\hskip 1pt}}
\def \I{{\rm I\kern -2.2pt I\hskip 1pt}}
\def \F{{\rm I\kern -2.2pt F\hskip 1pt}}
\def \D{{\rm I\kern -2.2pt D\hskip 1pt}}
\def \B{{\rm I\kern -2.2pt B\hskip 1pt}}




\begin{document}
\title{Parametric Curves and Applications in CAGD: Introduction and some 
Special Issues}
\author{ $\begin{array}{ccc} \mbox{T. Recio}\thanks{
Partially supported by CICyT PB 92/0498/C02/01 (Geometr\'{\i}a 
Real y Algoritmos), Esprit/Bra 6846 (Posso), and TIC-1026-CE} & & 
\mbox{J.R. Sendra \& J. Sendra}\thanks{Partially supported by 
Univ. Alcal\'a Proy. 030/95} \\ 
\mbox{Dpto. de Matem\'aticas} & & \mbox{Dpto. de Matem\'aticas} \\ 
\mbox{Universidad de Cantabria} & & \mbox{Universidad de Alcal\'a} \\ 
\mbox{Santander 39071 (Spain)} & & \mbox{Madrid 28871 (Spain)}   \\ 
\mbox{recio@matsun1.unican.es} & & \mbox{mtsendra@alcala.es} 
\end{array}$}
\date{}
\maketitle
%\vspace{0.5 cm} 
In this paper we briefly introduce the notion and main 
problems related to parametric curves, and we focus on two issues specially 
interesting in applications: the problem of computing real reparametrization 
of real parametric curves, and the rationality analysis and parametrization 
of offset curves of parametric curves.

Let $\K$ be a field of characteristic zero and ${\cal K}$ its algebraic 
clousure. An algebraic variety ${\cal V}$ in the affine $n$-space 
${\cal K}^{n}$ implicitely defined by the system of polynomial equations
\[ {\cal V}=\{ f_{1}(x_{1},\ldots,x_{n})=0,\ldots, 
f_{r}(x_{1},\ldots,x_{n})=0 \} \]
$f_{1},\ldots,f_{r}\in \K[x_{1},\ldots,x_{n}]$, is called parametric if 
there exist rational functions $p_{1},\ldots,p_{n} \in 
{\cal K}(t_{1},\ldots,t_{m})$ such that 
$f_{i}(p_{1}(t_{1},\ldots,t_{m}),\ldots,p_{n}(t_{1},\ldots,t_{m}))$
is identically zero for $i=1,\ldots,r$. Then, we say that $
(p_{1},\ldots,p_{n})$ is a rational  parametrization of ${\cal V}$. 
Every parametrization of ${\cal V}$ can also 
be seen as a map ${\cal P}$ from ${\cal K}^{m}$ to ${\cal K}^{n}$ such that 
for almost every parameter value $\bar{t}=(t_{1},\ldots,t_{m})\in 
{\cal K}^{m}$ (i.e. for all but a finite number of exceptions) 
${\cal P}(\bar{t})=
(p_{1}(\bar{t}),\ldots,p_{n}(\bar{t}))$ is a point of the variety 
${\cal V}$.

Here we deal with parametric curves; more particularly with parametric 
plane curves. Geometric questions related to parametric curves in higher 
dimension spaces can be studied considering birationally equivalent plane 
curves. Thus, we are dealing with varieties ${\cal C}$ in 
${\cal K}^{2}$ that can be  represented implicitaly by a polynomial 
$f(x,y)\in \K[x,y]$ 
\[ {\cal C}=\{ (a,b)\in {\cal K}^{2} \,/\, f(a,b)=0 \} \]
or parametricaly as ${\cal P}(t)=(p_{1}(t),p_{2}(t))$, where 
$p_{1},p_{2}\in {\cal K}(t)$. 

If one starts from the implicit representation of the variety the question 
of computing  rational parametric representations arises 
(parametrization problem), and the preliminary subproblem of deciding 
whether the curve can be parametrized has also to be considered 
(rationality problem). Parametric curves, or rational curves, are exactly  
the irreducible curves of genus zero. The genus is a birational invariant 
of the curve, that measures the deficient between the maximum allowable 
limit of singularities and the number of singularities of the curve, 
counted properly with multiplicities (if the curve has degree $d$ then 
$\frac{(d-1)(d-2)}{2}$ bounds the number of double points). Thus, 
parametric  curves  are the irreducible curves with the maximum allowable 
number of singularities; for instance: conics, cubics with one double 
point, $\ldots$. Hence, the rationality of the curve can be decided 
analyzing the multiplicity of the singularitites. Recently, several 
algorithms for computing the genus and rational parametrizations have 
been proposed [1],[13],[17],[21],[22]. 

If one starts from the parametric representation of the variety, the 
natural question of determining the  defining polynomial, or defining 
system of polynomials in the higher dimension case, appears 
(implicitization problem).  Implicitization procedures can be found for 
instance in [6],[11]. Also, many authors have addressed the problem of 
computing {\it good} parametrizations; where {\it good} affects to the 
degree of the rational functions in the parametrization, to the algebraic 
field extension where the parametrization is expressed, and to the  
coefficient lenght in the parametrization.  Proper parametrizations (i.e. 
parametrizations with optimal degree) are reached in [1],[13],[17],[21],
[22]. Reparametrization methods for improperly parametrized curves, and 
 algorithmic criteria for the existence of polynomial parametrizations are  
also available, for instance, in [2],[6],[11]. In [22] a procedure for 
determining parametrizations expressed over the smallest possible field 
extension of the ground field is presented. Similar theoretical results 
can be found in [17]. Nevertheless, purely arithmetic questions (for 
instance, the determination of upper bounds of the lenght of the 
coefficients of the rational functions involved in the parametrization) 
are still open questions. 

In addition, complexity analysis of some of these problems have been 
presented in [14],[23], and computer packages for dealing with algebraic 
curves and rational function fields are developed in [4],[10].

Parametric curves play an important role in the field of Computer Aided 
Geometric Design. They arise in practical applications as computer 
graphics, geometric modeling or manipulation of offset varieties. The 
natural field of applications of rational curves is Computer Aided 
Design.  Indeed, real varietes are the most common ones in Geometric 
Design. However, since Algebraic Geometry is mainly developed over 
algebraically closed fields, the majority of the previous 
algorithms are checked only over  the complex numbers. Thus, the usually 
forgotten reality issues are not treated. Special exceptions are: 
[3],[5],[9] where the existence of real parametrizations is analyzed, 
[16] where a real reparametrization algorithm is presented, 
and [22] where a  parametrization algorithm, that computes real 
parametrizations when they exist, is given. 

An interesting problem to mention is the algebraic manipulation of offset  
curves. Offset curves of a given curve, or more formerly parallel curves, 
are defined as the envelopes of systems of circles with centres on the 
original curve and fixed radious. They  arise in practical applications 
as tolerance a\-na\-ly\-sis, geometric control, robot path-planning and 
numerical-control machining problems.  In [7],[8] the principal 
geometric and topological properties of offset curves are stated, and  
an algorithmic approach, based on resultants, to compute the implicit 
equation of an offset plane curve is provided. Similarly, in [12] it is 
shown how to use  Computer Algebra techniques as Gr\"obner Basis (or any 
other algebraic equation solver as Characteristic Sets)  to determine 
implicit equations of the offset of rational surfaces. One of the main 
problems in offset curves manipulation is that the rationality of the 
original curve is not preserved (in general) when the offset curve is 
considered. Pottman gives explicit formulas to produce rational plane 
curves and surfaces whose offsets have rational components [15]. 
Nevertheless, it is not algorithmically clear how to deduce if a given 
curve has an offset with  rational components, and how to compute a 
rational offset parametrization, if they exist. In [19],[20] a complete 
rationality analysis of the offset curves is presented, and symbolic 
algorithms to parametrize the rational components of the offset curves 
to parametric curves are derived. 

In the sequel, we present the main ideas of the real reparametrization 
algorithm developed in [16], and we briefly summarize the rationality 
analysis and parametrization algorithm for offset curves given in [19],[20].
\newpage
%\vspace{1 mm}
%\newline
\begin{center}
{\bf Real Reparametrizations of Curves}
\end{center}
\vspace{1 mm}
Let ${\cal C}$ be a plane curve 
defined by a real polynomial $f \in \R[x,y]$, with $f$ irreducible 
over the complexes. Let us assume that ${\cal C}$ is rational, 
and let ${\cal P}(t)$ be a rational  proper complex 
parametrization of ${\cal C}$. Then, we study the existence 
and computation --by re\-pa\-ra\-me\-tri\-za\-tions of ${\cal P}(t)$-- 
of real parametrizations of ${\cal C}$.  For  this purpose, we associate 
with ${\cal C}$ an additional curve  that provides the complex parameter 
values that generates --via ${\cal P}$-- the real points on ${\cal C}$. 

More precisely, let $\bar{t}=t_1 + i\,t_{2},\,\, t_1,t_2 \in \R$, denote a 
generic complex number. Then the parametrization ${\cal P}(t)$ can be 
written in the form:
\[ {\cal P}(\tilde{t})=(\frac{h_{1}(t_1,t_2)+i\,\, 
f_{1}(t_1,t_2)}{m_{1}(t_1,t_2)}, \frac{h_{2}(t_1,t_2)+i\,\, 
f_{2}(t_1,t_2)}{m_{2}(t_1,t_2)}) \]
where $m_1,m_2,h_1,h_2,f_1,f_2 \in \R[x,y]$. Now, if ${\cal C}$ is 
parametrizable over the reals, there exist infinitely many  points 
$(t_1,t_2)\in \R^{2}$ such that ${\cal P}(t_1+i\,t_{2})$ is a real point 
on ${\cal C}$. Therefore 
the curves ${\cal F}_{1}$, ${\cal F}_{2}$, defined respectively by 
$f_{1}(t_1,t_2)$ and $f_{2}(t_1,t_2)$, have infinitely many common 
points, and hence, they have common components. Let  $\tilde{{\cal C}}$ 
be the curve defined by the common components of ${\cal F}_{1}$ and 
${\cal F}_{2}$. Observe that $\tilde{{\cal C}}$ has  real points, and that 
the mapping 
\[ \begin{array}{ccc}
\varphi:\tilde{{\cal C}} & \longrightarrow & {\cal C} \\ 
\,\,\,\,\,\,\,\,\,\,\,\,(a,b) &\longrightarrow & {\cal P}(a+ \,i \,b) 
\end{array} 
\]
is a one-to-one mapping  of the set of almost all  real points of 
$\tilde{{\cal C}}$ (i.e. all but a finite number of exceptions) onto 
the set of almost all  real points of ${\cal C}$.

The following theorem states the algebraic properties of 
$\tilde{{\cal C}}$ and the mapping $\varphi$ that allow us to construct 
the real reparametrization algorithm. 
\vspace{1 mm}
\newline
{\sc Theorem 1.} 
\vspace{1 mm}
\newline
\hspace*{8 mm}
(1) $\tilde{{\cal C}}$ has exactly one real component. 
Furthermore, it is either a real circle or a real line. In the 
\newline
\hspace*{13 mm} sequel, although it may be a line,  we will refer to the 
real component ${\cal C}^{\star}$ of $\tilde{{\cal C}}$ as the 
{\it associated} 
\newline
\hspace*{13 mm} {\it real conic} with ${\cal C}$. 
\vspace{1 mm}
\newline
\hspace*{8 mm}
(2)   ${\cal C}$ is parametrizable over $\R$ if and only if 
$\tilde{{\cal C}}$ has a real circle or a real line as a component.
\vspace{1 mm}
\newline
\hspace*{8 mm}
(3)  If ${\cal C}$ is parametrizable over $\R$ then for every 
real parametrization 
$(\frac{m_{1}(t)}{m_{2}(t)},\frac{m_{3}(t)}{m_{4}(t)})$ of 
${\cal C}^{\star}$ it 
\newline
\hspace*{13 mm}
holds that ${\cal P}(\frac{m_1}{m_2}+
i\,\frac{m_3}{m_4})$ is a proper real parametrization of ${\cal C}$.
\vspace{2 mm}
\newline
Thus, since ${\cal C}^{*}$ is rational (it is a circle or a line) and it 
is easy to check whether it has real points. The previous ideas can be 
summarized in an algorithm. Given a proper complex parametrization 
${\cal P}(t)$ of an affine plane curve ${\cal C}$, the algorithm  
decides whether ${\cal C}$ can be parametrized over the reals, and if so a 
proper real parametrization is computed.
\newline
\underline{Algorithm} 
\newline
1.  Take $\bar{t}=t_1+\,i\,t_2$ and compute
${\cal P}(\bar{t})=(\frac{h_{1}(t_1,t_2)+i\,\, 
f_{1}(t_1,t_2)}{m_{1}(t_1,t_2)}, \frac{h_{2}(t_1,t_2)+i\,\, 
f_{2}(t_1,t_2)}{m_{2}(t_1,t_2)})$ 
\newline
\hspace*{3 mm} 
where $m_1,m_2,h_1,h_2,f_1,f_2 \in \R[x,y]$.
\newline
2. Find the real factors of  $\tilde{g}(t_1,t_2)=gcd(f_1,f_2)$. 
\newline
3. If $\tilde{g}$ has neither a real circle  nor a line as a  real factor, 
 then {\sc Return} that ${\cal C}$ is not real. 
\newline
4. Take $g^{\star}(t_1,t_2)$ as the real circle factor or the real line 
 factor of $\tilde{g}(t_{1},t_{2})$. 
\newline
5. Obtain a real proper parametrization 
$(\frac{m_{1}(t)}{m_{2}(t)},\frac{m_{3}(t)}{m_{4}(t)})$   
 of   $g^{\star}(t_1,t_2)$. 
\newline
7. {\sc Return} ${\cal P}(\frac{m_1}{m_2}+
i\,\frac{m_3}{m_4})$.
\newpage
%\vspace{1 mm}
%\newline
\begin{center}
{\bf Parametrization  of Offset Curves}
\end{center}
\vspace{1 mm}
Let ${\cal C}$ be a parametric curve over ${\cal K}$ given by ${\cal P}(t)=
 (p_{1}(t),p_{2}(t)), p_{1},p_{2} \in {\cal K}(t)$ (we exclude the trivial 
case where ${\cal C}$ is a line). Then, the offset 
curve to ${\cal C}$ at distance $d$ is the set of all 
points of the affine plane  over ${\cal K}$ given by the formula 
\[ {\cal P}(t) \pm \frac{d}{m(t)} \,{\cal N}(t),
\]
where ${\cal N}(t)=(-p'_{2}(t), p'_{1}(t))$ is the normal vector,  and 
$m(t)= \sqrt{p'_{1}(t)^{2}+p'_{2}(t)^{2}}$ is its length. In the sequel, we 
denote by ${\cal O}_{d}$ the offset curve to ${\cal C}$ at distance $d$. 
Note that ${\cal O}_{d}$ is the envelope of the system of circles with 
centre on ${\cal C}$ and  radious $d$. 

One of the main problems in offset curves manipulation is that the 
rationality of the original curve is not preserved (in general) when the 
offset curve is considered. In the following we give a characterization of 
those parametric curves with parametric offsets, as well as an algorithm 
for parametrizing parametric offsets. 

For this purpose,  we introduce the 
following notion: let ${\cal P}(t)=(p_{1}(t),p_{2}(t))$ be a proper 
parametrization of ${\cal C}$. Then, {\it the curve of reparametrization of 
${\cal O}_{d}$  associated with ${\cal P}$} is defined as the plane curve 
given by the primitive part --w.r.t $y$-- of the  numerator of 
$2p_{2}'(x)y+p_{1}'(x)(y^{2}-1)$, and it is denoted as 
${\cal G}_{{\cal P}}$. Then one has the following characterizations:
\vspace{1 mm}
\newline
{\sc Theorem 2.}  The following statements are equivalent 
\vspace{1 mm}
\newline
\hspace*{8 mm} 
(1) All the components of ${\cal O}_{d}$ are rational 
\vspace{1 mm}
\newline
\hspace*{8 mm} 
(2) There exists  $\varphi \in {\cal K}(t)$ of degree at most two, 
and a rational proper parametrization ${\cal P}(t)$ of 
\newline
\hspace*{13 mm} 
${\cal C}$ such that the normal vector of ${\cal P}(\varphi(t))$ has 
rational length.
\vspace{1 mm}
\newline
\hspace*{8 mm} 
(3) There exists a  proper parametrization ${\cal P}$  of ${\cal C}$ 
such that  at least one component of ${\cal G}_{{\cal P}}$ is 
\newline
\hspace*{13 mm} rational. Furthermore, if $(R_{1},R_{2})$ parametrizes 
the component of ${\cal G}_{{\cal P}}$ then ${\cal P}(R_{1}(t))$ has 
\newline
\hspace*{13 mm} normal vector with rational length.
\vspace{2 mm}
\newline
As a consequence of theorem 2  one may deduce that the rationality of the 
offset curves does not depened on the distance. That is, either all the 
components of ${\cal O}_{d}$ are rational at any  nonzero distance,  or  
both components have positive genus  at any  nonzero distance. Therefore, 
the geometric analysis of the offest can be done at a fixed distance. 
Indeed, one may be more precise and distinguishes the cases of rational 
offsets or offsets with two rational components.
\vspace{1 mm}
\newline
{\sc Theorem 3.} The following statements are equivalent
\vspace{1 mm}
\newline
\hspace*{8 mm}
(1)  ${\cal O}_{d}$ is rational 
\vspace{1 mm}
\newline
\hspace*{8 mm}
(2) There exists a rational parametrization  ${\cal P}$ of ${\cal C}$
such that ${\cal G}_{\cal P}$ is rational.
\vspace{1 mm}
\newline
\hspace*{8 mm}
(3) For every rational parametrization ${\cal P}$ of ${\cal C}$ the curve 
${\cal G}_{\cal P}$ is rational.
\vspace{2 mm}
\newline
{\sc Theorem 4.} The following statements are equivalent:
\newline
\hspace*{8 mm}
(1)  ${\cal O}_{d}$ has two rational components
\vspace{1 mm}
\newline
\hspace*{8 mm}
(2) There exists a rational proper parametrization  ${\cal P}$ of 
${\cal C}$ such that ${\cal G}_{\cal P}$ is reducible.
\vspace{1 mm}
\newline
\hspace*{8 mm}
(3) For every rational proper parametrization ${\cal P}$ of ${\cal C}$ 
the curve ${\cal G}_{\cal P}$ is reducible.
\vspace{2 mm}
\newline
We finish this paper  with the algorithm  that summarizes the previous 
ideas. Given a proper  parametrization ${\cal P}(t)$ of an affine plane 
curve ${\cal C}$, it decides whether the components of ${\cal O}_{d}$  
are rational, and if so,  rational  parametrizations of all the components 
are computed.
\newline
\underline{Algorithm} 
\footnotesize
\begin{tabbing}
aa \= 1. \=  aaaa \= 5.3. \= aaaa \= \kill
\\
\> 1. \> ${\cal N}:=(-p'_{2},p'_{1})$; {\bf If} $|| {\cal N} || \in 
{\cal K}(t)$ {\bf then return} \\
\> \> \>  $\ll$ ${\cal O}_{d}$ has two rational components; 
${\cal P}\pm \frac{d}{|| {\cal N} ||} {\cal N}$ parametrizes the 
components of ${\cal O}_{d}$ $ \gg$   \\
\> 2. \> Compute the genus  $g$ of the curve ${\cal G}_{\cal P}$
defined by  $2p'_{2}(x)y+p'_{1}(x)(y^2-1)$. \\
\> 3. \>   {\bf If} $g\neq 0$ {\bf then return} $\ll$ no component of 
${\cal O}_{d}$ is rational $\gg$  {\bf else} \\
\> \> \> 3.1. \> Determine a rational parametrization 
${\cal R}=(R_{1},R_{2})$ of  ${\cal G}_{\cal P}$. \\
\> \> \> 3.2. \> {\bf Return} $\ll {\cal O}_{d}$ is rational and  
${\cal Q}={\cal P}(R_{1})+ \frac{2\,d\,R_{2}}{n_{2}(R_{1})(R^2_{2}
+1)}{\cal N}(R_{1})$ parametrizes ${\cal O}_{d}$ $\gg$.
\end{tabbing}
\footnotesize
\begin{thebibliography}{20}
\bibitem{} Abhyankar S.S, Bajaj C.L., (1988), {\it Automatic 
parametrization of rational curves and surfaces III: Algebraic plane 
curves}. Computer Aided Geometric Design {\bf 5}, 390-321.
\bibitem{} Abhyankar S., (1990), {\it Algebraic Geometry for scientists 
and engineers}. Mathematical surveys and monographs, 35. A.M.S. 1990.
\bibitem{} Alonso C., Guti\'errez J., Recio T., 
{\it Reconsidering algorithms for real parametric curves}. To appear  
in J. AAECC.
\bibitem{} Alonso C., Guti\'errez J., Recio T., (1994), {\it FRAC: 
A Maple package for computing in the rational function field $K(X)$}. 
MapleV: Mathematics and Its Applications, pp. 107-115. (Ed. R.J. Lopez)  
Birkhauser (1994).
\bibitem{} Bajaj C., Royappa A. V., (1992), {\it Finite 
representations of real parametric curves and surfaces}. Technical 
Report, CAPO Report CER-92-28. Purdue University, september 1992.
\bibitem{} Cox D., Little J., O'Shea D., (1992), {\it Ideals, 
varieties and algorithms (An Introduction to Computational Algebraic 
Geometry and Commutative Algebra)}. Springer Verlag
\bibitem{} Farouki R. T., Neff C.A. (1990), {\it Analytic properties of 
plane offset curves}. Comput. Aided Geom. Design {\bf 7} 83-99.
\bibitem{} Farouki R. T., Neff C.A. (1990), {\it Algebraic properties of 
plane offset curves}. Comput. Aided Geom. Design {\bf 7} 100-127.
\bibitem{} Guti\'errez J., Recio T., (1992), {\it Rational 
function decomposition and Groebner basis in the parametrization 
of a plane curve}. LATIN 92. L.N.C.S. 583, Springer-Verlag. 1992, 
pp. 231-246.
\bibitem{} Gebauer R., Kalkbrener M., Wall B. and Winkler F.,(1991), 
{\it CASA: A Computer Algebra Package for Constructive Algebraic 
Geometry}. In S. M. Watt  (ed.) Proc. ISSAC'91, 403-410, ACM Press.
\bibitem{} Hoffmann C., (1989) {\it Geometric and Solid Modeling: An 
Introduction}. Morgan Kaufmann, USA (1989). 
\bibitem{} Hoffman C. (1990), {\it Algebraic and Numerical Techniques for 
Offsets and Blends}. W. Dahmen et al. (eds.), Computation of Curves and 
Surfaces, (Kluwer Ac. Publ.) 499-528.
\bibitem{} van Hoeij M. (1994), {\it Computing parametrizations of 
rational algebraic curves}. In J. von zur Gathen (ed.), Proc. ISSAC94 
187-190, ACM Press.
\bibitem{} M\"nuk M., Sendra J.R., Winkler F. (1993), {\it On the 
Complexity of Parametrizing Curves}. Techn. Rep. RISC 93-51, Research 
Inst. Symb. Comp., Univ. Linz.
\bibitem{} Pottman H.,  (1994), {\it Rational curves and surfaces with 
rational offsets}. Tech. Rep. Nr.1. Institut f\"ur Geometrie, Technische. 
Univ. Wien.
\bibitem{} Recio T., J.R. Sendra (1995) {\it Real Reparametrizations of 
Real Curves}. Technical Report. Dep. Matem\'aticas, Universidad de 
Cantabria.
\bibitem{} Schicho J., (1992) {\it On the choice of pencils in 
the parametrization of curves}. J. Symbolic Computation 
{\bf  14,} 557-576.
\bibitem{} Sederberg, T.W. (1986), {\it Improperly parametrized 
rational curves}. Computer Aided Geometric Design {\bf 3}, 67-75.
\bibitem{} Sendra J.R., Sendra J. (1995), {\it On the  Rationality of 
Offset  Curves}. Techn. Rep. RISC 95-02, Research 
Inst. Symb. Comp., Univ. Linz.
\bibitem{} Sendra J.R., Sendra J. (1995), {\it Curvas Offset Racionales}. 
Preprint.
\bibitem{} Sendra J.R., Winkler F. (1991), 
{\it Symbolic Parametrization of Curves}. J. Symbolic Computation 
{\bf 12/ 6,} 607-631.
\bibitem{} Sendra J.R., Winkler F. (1994), 
{\it Optimal Parametrization of Algebraic Curves.} 
Techn. Rep. RISC 94-65, Research Inst. Symb. Comp., Univ. Linz.
\bibitem{} Teitelbaum J. (1989), {\it On the Computational Complexity of 
the resolution of plane curve singularities}. ISSAC-89 LNCS no. 358, 
pp. 285-292. Springer Verlag
\bibitem{} Walker R.J. (1950), {\it  Algebraic Curves.} Princeton 
Unversity Press.
\end{thebibliography}
\end{document}
